\scanpage{12}
\mspage{1)}
Stratified spaces (study by \deleted{tubular}
\inserted{conical} neighborhoods)
\textcircled{I}

\textcircled{1} Let $Y$ be a topological space. A \inserted{\uline{\unsure{conical}}} \uline{\deleted{tubular} neighborhood}
\inserted{(``\uline{absolute}'')} of \deleted{$Y$}
\inserted{(or ``cone over $Y$'')} is a topological space
$T\deleted{\unsure{_Y}}$ containing $Y$,
(the ``\uline{tube}''), equipped with a closed subset $\dot T$
\deleted{\illegible{1}} (the ``boundary'' of the tube'') and an action of the topological monoid
$\mathbf{I}=[0,1]\,\unsure{}$ (under multiplication) on $T$, denoted
% P08: present overworked mark after [0,1]; identity and cancellation state uncertain.
\[
(1)\qquad (\lambda,t)\longmapsto\lambda t
\qquad \mathbf{I}\times T\longrightarrow T
\]
satisfying :
\[
\left.\begin{array}{ll}
a)&\lambda y=y\quad\forall y\in Y\\
b)&0.\deleted{\unsure{T|\bar T}}T\subset Y
\end{array}\right\}
\]
\margin{thus the $\lambda_T:T\longrightarrow T$
$(0\leq\lambda\leq1)$ form a homotopy from
the identity $\operatorname{id}_T$ (for $\lambda=1$)
to a \uline{retraction} of $T$ onto $Y$,
leaving the points of $Y$ fixed during the homotopy}
\margin{\uline{NB} Condition b) in \unsure{its|the} weak form
$0.\dot T\subset Y$, plus condition c)
(to which this gives a meaning) suffice :
imply}

c) The \deleted{map} diagram of maps
\diagram{\[
\unsure{(2)|(b)}\qquad
\begin{tikzcd}[ampersand replacement=\&,column sep=large,
execute at end picture={
\path (p012a.east) -- coordinate[pos=.32] (p012extrahead) (p012b.west);
\draw ([xshift=-3pt,yshift=2pt]p012extrahead) -- (p012extrahead) -- ([xshift=-3pt,yshift=-2pt]p012extrahead);
}]
|[alias=p012a]| \mathbf{I}\times\dot T \arrow[r,"{(\lambda,\unsure{x})\mapsto\lambda x}",
"\deleted{\unsure{g|f}}"'] \& |[alias=p012b]| T\\
\{0\}\times\dot T \arrow[u,<->] \arrow[r,"{(0,x)\mapsto0.x}"'] \& Y \arrow[u,hook,"\mathrm{inc}"']
\end{tikzcd}
\]}
is \uline{cocartesian}$^{*}$, i.e. identifies $T$ with the cone of the map
\[
(3)\qquad f:\dot T\longrightarrow Y\qquad f(\unsure{a|x})=0.x
\]
\margin{\uline{$^{*}$} This means that
$\mathbf{I}_{]0,1]}\times\dot T\xrightarrow{\substack{\deleted{\unsure{homeomorphism}}\\\mathrm{bijective}\\\sim}}T\setminus Y$
and that a subset of $T$ is \unsure{open|open} if and only if its
\deleted{\unsure{its intersection|intersection}} \unsure{inverse|inverse}
\inserted{\unsure{images}} in $\mathbf{I}\times\dot T$ and in $Y$ are}
\[
(4)\qquad T\simeq C(f)
\]
Under this identification, $Y\hookrightarrow\deleted{\unsure{T}}T
\hookleftarrow\dot T$ correspond
to the canonical maps \deleted{\unsure{of the cone $f$}}
$\operatorname{but}(f)\hookrightarrow C(f)\hookleftarrow
\operatorname{source}f$
and the action of $\mathbf{I}$ on $T$ comes from the action of
$\mathbf{I}$ on $C(f)$, defined by passage to the quotient of
the obvious action of $\mathbf{I}$ on $\mathbf{I}\times\dot T$, and the trivial action
of $\mathbf{I}$ on $Y$. The axioms a) b) c) also mean
that the situation $(Y\mathrel{\tikz[baseline=-.45ex]{\draw (.12,.09) .. controls (-.01,.09) and (-.01,0) .. (.12,0) -- (.65,0);}}T
\mathrel{\tikz[baseline=-.45ex]{\draw (0,0) -- (.53,0) .. controls (.68,0) and (.68,.11) .. (.53,.13);}}\dot T,
\mathbf{I}\times T-T)$
is isomorphic : that of a $C(f)$ with $f:\dot T-Y$.
Note that $\dot T$ is known when one knows the action
of $\mathbf{I}$ on $T$, \deleted{\unsure{as one}} on setting
\[
(5)\qquad \mathring T=T\setminus\dot T
\quad\text{(the \uline{interior} of the tube)}
\qquad\text{and note that }\mathring T=\bigcup_{\lambda\in[0,1[}\lambda T
\]
\margin{11}
